\input{preamble.tex}

\begin{document}

% Indicate your name and your student number here (e.g., A01xxx)
\header{Assignment 1}{Wira Azmoon Ahmad}{A0149286R}

%======================================
\section{Pure Strategy Nash Equlibria}

\subsection{}
$T,L$ with payoff $7,1$

\subsection{}
$T,L$ with payoff $5,5$ \\
$B,R$ with payoff $4,4$

\subsection{}
$T,A$ with payoff $5,5$

%======================================
\newpage
\section{Ultimatum Game}

\subsection{Players, Actions and Utilities}
\subsubsection{Players}
\[N=\{proposer (P), responder (R)\}\]
\subsubsection{Actions}
\[A_P=\{0,...,10\}\]
\[A_R=\{accept, reject\}\]
\subsubsection{Utilities}
\[u_P(i,x) =
    \begin{cases}
        i, & \text{if } x=accept\\
        0, & \text{otherwise}
    \end{cases}
\]
\[u_R(i,x) =
    \begin{cases}
        10-i, & \text{if } x=accept\\
        0, & \text{otherwise}
    \end{cases}
\]

\subsection{Pure Strategy Nash Equilibria}
Two exist:
\[\vec{a}_1=(10,accept),u_P(\vec{a}_1)=10,u_R(\vec{a}_1)=0\]
\[\vec{a}_2=(10,reject),u_P(\vec{a}_2)=u_R(\vec{a}_2)=0\]

\subsection{Behaviour}
No. The \textit{proposer} may want to choose $i$ such that the \textit{responder} benefits as well.
%======================================
\section{Pure and Mixed Nash Equilibria}
Let the row player pick T with probability $p$ and column player pick L with probability $q$.
We have 
\[\vec{p}_r=(p,1-p)\]
\[\vec{p}_c=(q,1-q)\]
\subsection{}



\subsubsection{Pure}
None.
\subsubsection{Mixed}
Compare the expected utility for the row player, 
of the column player picking each action.
\[u_r(\vec{p}_r,(1,0))=7p\]
\[u_r(\vec{p}_r,(0,1))=2p+3(1-p)=3-p\]
The two values must be equal so that the row player can be indifferent for their choice.
\[7p=3-p\]
\[\implies p=3/8\]
Similarly for the column player,
\[u_c((1,0),\vec{p}_c)=q+2(1-q)=2-q\]
\[u_c((0,1),\vec{p}_c)=5q+(1-q)=4q+1\]
\[2-q=4q+1\]
\[\implies q=1/5\]
Then, the mixed Nash equilibrium strategy profile is
\[\vec{p}=\begin{pmatrix}
    3/8 & 1/5\\
    5/8 & 4/5
\end{pmatrix}\]
with expected utilities
\[u_r(\vec{p})=7\left(\frac{3}{8}\right)=\frac{21}{8}\]
\[u_c(\vec{p})=2-\frac{1}{5}=\frac{9}{5}\]

\subsection{}
\subsubsection{Pure}
$T,L$ with payoff $3,5$ \\
$B,R$ with payoff $5,3$
\subsubsection{Mixed}
\[u_r(\vec{p}_r,(1,0))=3p+(1-p)=2p+1\]
\[u_r(\vec{p}_r,(0,1))=5(1-p)=5-5p\]
\[2p+1=5-5p\]
\[\implies p=4/7\]
\[u_c((1,0),\vec{p}_c)=5q+(1-q)=4q+1\]
\[u_c((0,1),\vec{p}_c)=3(1-q)=3-3q\]
\[4q+1=3-3q\]
\[\implies q=2/7\]
\[\vec{p}=\frac{1}{7}
\begin{pmatrix}
    4 & 2\\
    3 & 5
\end{pmatrix}\]
\[u_r(\vec{p})=2\left(\frac{4}{7}\right)+1=\frac{15}{7}\]
\[u_c(\vec{p})=4\left(\frac{2}{7}\right)+1=\frac{15}{7}\]

\subsection{}
\subsubsection{Pure}
$T,L$ with payoff $10,5$ \\
$C,R$ with payoff $3,7$
\subsubsection{Mixed}
Row player will not choose B since 
\[u_r(B,\vec{q})<0.5u_r(B,\vec{q})+0.5u_r(B,\vec{q})\]
Column player will not choose M since
\[u_c(\vec{p},M)<0.5u_c(\vec{p},L)+0.5u_c(\vec{p},R)\]

From the remaining choices,
\[u_r(\vec{p}_r,(1,0))=10p+(1-p)=9p+1\]
\[u_r(\vec{p}_r,(0,1))=2p+3(1-p)=3-p\]
\[9p+1=3-p\]
\[\implies p=1/5\]
\[u_c((1,0),\vec{p}_c)=5q+3(1-q)=2q+3\]
\[u_c((0,1),\vec{p}_c)=4q+7(1-q)=7-3q\]
\[2q+3=7-3q\]
\[\implies q=2/5\]
\[\vec{p}=\frac{1}{5}
\begin{pmatrix}
    1 & 2\\
    4 & 3
\end{pmatrix}\]
\[u_r(\vec{p})=3-\frac{1}{5}=\frac{14}{5}\]
\[u_c(\vec{p})=2\left(\frac{2}{5}\right)+3=\frac{19}{5}\]

%======================================
\section{Two Player Normal Form Game}

\end{document}